\section{Method}
\label{sec:method}

MouseClaimBench is a claim-level \ac{dss}. Its inputs are a candidate statement, a domain profile, and the artifacts that may support the statement. Its output is an auditable disposition with the contributing observations, rules, sources, and an editorial rationale. It does not calculate a paper-level acceptance score. Fig.~\ref{fig:claimbench-workflow} presents the workflow.

\begin{figure}[t!]
\centering
\fbox{%
\begin{minipage}{0.92\textwidth}
\small
\centering
\begin{tabularx}{0.90\textwidth}{>{\centering\arraybackslash}X c >{\centering\arraybackslash}X c >{\centering\arraybackslash}X}
\textbf{Manuscript and artifacts} & $\rightarrow$ &
\textbf{Candidate claims} & $\rightarrow$ &
\textbf{Evidence contracts} \\
\multicolumn{5}{c}{$\downarrow$} \\
\textbf{External controls} & $\rightarrow$ &
\textbf{Non-compensatory gate} & $\rightarrow$ &
\textbf{Decision artifact} \\
\end{tabularx}
\vspace{0.8em}

\begin{tabularx}{0.90\textwidth}{p{0.25\textwidth}X}
\textbf{Gate input} & Prediction, reproducibility, evidence retrieval, topology specificity, direction, causality, uncertainty, manuscript wording, cost, and release status. \\
\textbf{Gate output} & Supported, blocked, or uncertain claim status with a pointer to the artifact that justifies the decision. \\
\textbf{LLM boundary} & LLMs can suggest candidate claims, but they cannot authorize scientific interpretations. \\
\end{tabularx}
\end{minipage}}
\caption{ClaimBench decision-support workflow. The framework separates claim detection from claim authorization and turns manuscript-level interpretations into executable decision artifacts.}
\label{fig:claimbench-workflow}
\end{figure}


\subsection{Domain-aware evidence contracts}
\label{subsec:claim-contracts}

Let $c$ denote a candidate statement, $k(c)$ its claim family, and $d(c)$ the domain protocol under which its evidence was generated. A contract is defined as
\begin{equation}
\mathcal{C}(c)=\langle k,d,R_k,\Phi_d,w^{+},w^{-},\mathcal{A}\rangle,
\label{eq:claim-contract}
\end{equation}
where $R_k$ is the set of evidence blocks required by claim family $k$, $\Phi_d$ is the set of domain-specific predicates, $w^{+}$ and $w^{-}$ are admissible and blocked wording, and $\mathcal{A}$ contains the source artifacts and revisions. The contract is executable because every required block names both a source and a decision rule.

The observed value for block $j$ is denoted by $x_{j,d}$. It remains in the scale reported by its source analysis. The contract does not assume that predictive correlation, a cross-mouse coefficient, a matched structure--function delta, and a bootstrap interval are commensurable. Instead, the domain predicate produces one evidence state
\begin{equation}
B_{j,d}(c)=\phi_{j,d}\bigl(x_{j,d},\mathcal{A}_{j}\bigr)
\in \mathcal{S}_{E},
\qquad
\mathcal{S}_{E}=\{P,F,U,N,R\},
\label{eq:evidence-state}
\end{equation}
where $P$ means passed, $F$ failed, $U$ unknown, $N$ not applicable, and $R$ requires external review. A failed state means that the stated test was executed and did not pass. An unknown state means that a required observation is absent. A not-applicable state means that the protocol did not target the block. These states are not interchangeable.

The v3 taxonomy contains ten authorization targets. Prediction, topology specificity, direction, local structure--function association, mechanism, causality, and digital-twin scope describe scientific interpretation. Computational reproducibility, internal reproduction, and external replication describe distinct forms of repeatability and confirmation. Table~\ref{tab:claim-semantics} gives their required blocks. A mechanistic claim requires prediction, internal reproduction, topology specificity, and directed identifiability. A digital-twin claim also requires intervention, whole-brain coverage, independent validation, and reproducible computation.

\begin{table}[t!]
\centering
\caption{Author-proposed operational requirements in profile v1.1. Every listed block must pass for authorization under this profile. The table does not assert community consensus or universal scientific sufficiency.}
\label{tab:claim-semantics}
\footnotesize
\renewcommand{\arraystretch}{1.02}
\begin{tabularx}{\textwidth}{@{}>{\raggedright\arraybackslash}p{0.19\textwidth}YY@{}}
\toprule
Claim & Required evidence block & Boundary preserved by the contract \\
\midrule
Predictive & Held-out improvement & No topology, direction, or cause \\
Computationally reproducible & Clean revision and repeatable workflow & Execution is not identification \\
Within-resource reproduction & Non-overlapping units or cohorts in one resource & Subtypes and sampling units remain explicit. This is not external replication \\
Externally replicated & Independent resource, laboratory, or study & Hold-outs within one resource are insufficient \\
Topology specific & Prespecified topology controls & Association alone is insufficient \\
Directed & Direction-identification test & Temporal order alone is insufficient \\
Local observational structure--function & Matched controls and dependence-aware stability & Local association remains non-causal \\
Mechanistic & Prediction, within-resource reproduction, topology, direction & Any missing block vetoes support \\
Causal & Intervention or causal-identification design & Correlation cannot authorize causality \\
Complete entity-specific digital twin & Prediction, mechanism, intervention, whole-brain coverage, entity specificity, independent validation, reproducible compute, and operational compute & A partial cortex model or generic archetype is not a complete twin \\
\bottomrule
\end{tabularx}
\end{table}


The separation among reproducibility concepts is deliberate. A clean revision supports computational reproducibility. Consistency across non-overlapping mice or windows in one public resource supports internal reproduction. Replication requires a different resource, laboratory, or study. A result may pass the first two blocks while the third remains unknown. This is the situation for the available \ac{microns} hold-outs.

\subsection{Non-compensatory disposition policy}
\label{subsec:claim-authorization}

For claim family $k$, let $\mathbf{B}_{k}(c)=\{B_{j,d}(c):j\in R_k\}$ contain the states of all required blocks. The policy is
\begin{equation}
D_k(c)=
\begin{cases}
\text{blocked}, & F\in\mathbf{B}_{k},\\
\text{external review}, & F\notin\mathbf{B}_{k}\text{ and }R\in\mathbf{B}_{k},\\
\text{uncertain}, & F,R\notin\mathbf{B}_{k}\text{ and }U\in\mathbf{B}_{k},\\
\text{out of scope}, & F,R,U\notin\mathbf{B}_{k}\text{ and }N\in\mathbf{B}_{k},\\
\text{supported}, & \mathbf{B}_{k}=\{P\}.
\end{cases}
\label{eq:five-state-gate}
\end{equation}
The first line makes the policy non-compensatory. Once a required executed test fails, a larger value in another block cannot recover the claim. The ordering also preserves a useful distinction. A failed topology test blocks a mechanistic statement even when direction was not measured, because the statement is already unsupported. In the absence of failure, missing evidence remains uncertain rather than being converted into a negative measurement.

The policy is related to a multicriteria veto but differs from an outranking score \citep{figueira2013electre}. It does not compare alternatives and it contains no trade-off weights. A stronger claim simply has a larger set of mandatory blocks. The recommended manuscript action follows from the disposition. Supported wording remains bounded to the declared scope. Blocked wording is removed or weakened. Uncertain wording requires new evidence or explicit qualification. Out-of-scope wording is excluded from the evaluated contribution. External-review wording remains with a qualified human decision maker.

This workflow model is distinct from the legacy v2 numerical-stability analysis. The legacy analysis creates 11 samples from each constructed case. It evaluates the nominal vector and positive and negative perturbations of 0.03 for each of five numerical fields, changed one at a time. A claim is numerically stable-supported when at least 95\% of those samples pass, stable-blocked when at most 5\% pass, and locally uncertain otherwise. This three-state output is a deterministic neighbourhood diagnostic. It is not a posterior probability, confidence interval, conformal guarantee, or substitute for the five dispositions in Equation~\ref{eq:five-state-gate}.

Evaluation uses ordinary classification measures. For a reference set of supportable claims, a \ac{fp} is an unsupported claim authorized by a policy and a \ac{fn} is a supportable claim not authorized. The corresponding rates are
\begin{align}
\mathrm{FPR} &= \frac{FP}{FP+TN},
\label{eq:fpr}\\
\mathrm{FNR} &= \frac{FN}{FN+TP}.
\label{eq:fnr}
\end{align}
These are standard error rates. Earlier artifacts also called them the Overclaiming Risk Index and Conservativeness Index. Version 3 retains those keys only for backward compatibility and does not present them as new metrics.

\subsection{Artifact binding and authority boundaries}
\label{subsec:artifact-binding}

Evidence blocks are immutable records containing a name, state, source path, rule, rationale, and named observations. Duplicate block names are rejected. Missing required blocks are represented as unknown. The real-case builder requires a complete inventory and raises an error when an expected source artifact is absent. This design prevents silent defaults from becoming evidence.

Claim discovery precedes authorization. The reproducible mode uses deterministic patterns to identify candidate sentences associated with prediction, mechanism, causality, structure--function relations, replication, digital twins, and state-of-the-art language. An optional local \ac{json} file may contain candidates generated by an external \ac{llm}. Candidate source is recorded, but every candidate is untrusted. It cannot modify predicates, create observations, change provenance, or authorize a claim.

Human authority remains necessary at three points. Domain experts define whether a rule is scientifically appropriate. Adjudicators establish reference labels for a future human evaluation. Authors and reviewers retain responsibility for the final wording and significance of a study. The automated system can expose a missing intervention or failed topology control. It cannot decide novelty, ethical acceptability, practical importance, or manuscript acceptance.

The release gate enforces the same boundary. It verifies required artifacts, expected decisions, recorded revisions, and the absence of `-dirty' suffixes. It also retains active scientific blockers in the release output. A package can therefore be technically ready while human decision benefit and external biological replication remain unvalidated. This result is intentional. Technical readiness and scientific sufficiency are different decisions.
